% =====================================================================
% 04_faddeev_ags.tex
% 第 4 章 Faddeev 分解、AGS 方程与 U-算符
% Wave 1 / Agent B
% =====================================================================

\chapter{Faddeev 分解、AGS 方程与 U-算符}
\label{ch:04_faddeev_ags}

% =========================================
\section{物理动机}
\label{sec:ch04-motivation}

\paragraph{为什么三体 LS 方程不直接可解。}
\cref{ch:01_lippmann_schwinger} 已经在两体散射中建立了 Lippmann–Schwinger
（LS）方程的对应代数：将 T-算符 \(\Top\) 与势 \(\Vop\) 由
\(\Top = \Vop + \Vop\,\Gzero\,\Top\) 关联。把同一形式直接套到三体系统、
取 \(\Vop \equiv \Vij{12}{} + \Vij{23}{} + \Vij{31}{}\)（三对两体势之和），
得
\begin{equation}
\label{eq:ch04-LS-3body-naive}
\Top = \Vop + \Vop\,\Gzero\,\Top.
\end{equation}
这一方程在算子层面看起来与两体 LS 一模一样，但它\textbf{在动量表象
下是不可解（不数值化）的}。
原因是其核 \(\Vop\,\Gzero\) 含
\(\Vij{ij}{}\,\Gzero\) 这种"两体对在三体 Hilbert 空间内的嵌入"，每一项
\(\Vij{ij}{}\) 都包含一个 Dirac \(\delta\) 函数（"剩下那一个粒子的动量"
不发生散射）。具体写出动量积分核 \(K(\bvec{p}, \bvec{q}; \bvec{p}', \bvec{q}')\)：
\(\Vij{12}{}\) 项里 \(\bvec{q}\) 完全不参与积分变量，给出
\(\delta(\bvec{q} - \bvec{q}')\)；同理 \(\Vij{23}{}, \Vij{31}{}\) 各自给出
另一组 \(\delta\)。这种形如 "\(\delta\) 函数 + 紧 kernel"的混合算子
在 Banach 空间不是紧的（compact），因此 Fredholm 选择定理无法保证
\((\identity - \Vop \Gzero)\) 可逆。

这就是 \emph{disconnected kernel}（不连通核）问题：
LS 方程在三体下因核里的 \(\delta\) 函数失效。
Faddeev 1960、Lovelace 1964、Alt-Grassberger-Sandhas（AGS）1967 的工作把这一困境从根本上绕过。

\paragraph{Faddeev–AGS 的两个关键思想。}
\begin{enumerate}
  \item \textbf{Faddeev 分解}：把全波函数拆成三段，
        \(\psi = \psifad{1} + \psifad{2} + \psifad{3}\)，
        让每一段 \(\psifad{i}\) 满足只含一个两体子势 \(\Vij{jk}{}\)
        （\(i \neq j, k\)）的简单方程。三段之间通过\textbf{置换 \(\Pop\)}
        关联。这样原本"不连通的 \(\delta\) 函数"被转化成"通道之间的耦合"，
        kernel 变为紧算子。
  \item \textbf{AGS 整合}：把 Faddeev 三式重新组合成一组对称的
        通道间散射振幅 \(\Uampl{\beta}{\alpha}\)，再利用置换对称性把
        三式压成单一方程。这是最终被代码数值化的形式。
\end{enumerate}

\paragraph{本章的目标。}
本章把 Faddeev 与 AGS 推到求解器代码可直接消费的"\(U\)-方程"形式，并解释
求解器函数 \texttt{solve\_faddeev\_equations} 的签名为什么长成那个样子。
读完本章你应能：
\begin{itemize}
  \item 在纸笔上独立写出 Faddeev 三式与 AGS 形式之间的等价关系；
  \item 看到 \texttt{solve\_faddeev.h} 的接口立刻知道每一个参数对应数学
        上的哪个量（特别是 \texttt{P123\_sparse\_*}, \texttt{V\_WP\_*},
        \texttt{C\_WP\_*}, \texttt{q\_com\_idx\_array}）；
  \item 知道弹性通道 \(U_{0\beta}\) 与断裂通道 \(U_{0(0)}\) 的差别在
        哪条公式与哪个数组上；
  \item 知道 Tic-tac 求解器在 LS-form vs 算符 form 上的取舍。
\end{itemize}

% =========================================
\section{数学推导}
\label{sec:ch04-math}

% -----------------------------------------
\subsection{两体子 T 算符 \(\Tsub{i}\)}
\label{subsec:ch04-tsub}

定义"嵌入到三体 Hilbert 空间的两体 T 算符"
\begin{equation}
\label{eq:ch04-tsub-def}
\Tsub{i}(\zeplus) \;\equiv\; \Vij{jk}{} + \Vij{jk}{}\,\Gzero(\zeplus)\,\Tsub{i}(\zeplus),
\end{equation}
其中 \((i, j, k)\) 取 \((1,2,3)\) 的循环置换，\(\Vij{jk}{}\) 是
"对 \((j,k)\) 的两体势"（\(\bvec{q}^{(i)}\) 维度上是单位算符），
\(\Gzero(\zeplus) = (\zeplus - \Hzero)^{-1}\) 是三体自由 Green 算符，
\(\zeplus = E + i\eps\)。
\(\Tsub{i}\) 与两体 \(\Top_{jk}\) 的关系为
\(\Tsub{i} = \Top_{jk} \otimes \identity_{\bvec{q}^{(i)}}\)：
"在 \(\bvec{p}^{(i)}\) 上是真实的两体 \(T\)，在 \(\bvec{q}^{(i)}\) 上是恒等"。
代码里 \(\Tsub{i}\) 不直接显式出现——它被吸收到\textbf{Sturmian-Wave-Packet
(SWP) 系数} \(C_{\PWchan}\) 与势矩阵 \(V_{\PWchan, \PWchanp}\) 的乘积里
（即所谓 "VC-product"，详见 \cref{ch:11_faddeev_solver}）。

% -----------------------------------------
\subsection{Faddeev 三分量分解}
\label{subsec:ch04-faddeev}

设全 Hamilton 量 \(\Hop = \Hzero + \Vop = \Hzero + \sum_{(jk)} \Vij{jk}{}\)。
Lippmann–Schwinger 形式的全波函数
\begin{equation}
\label{eq:ch04-LS-full}
\ket{\Psi^{(+)}}
= \ket{\Phi} + \Gzero(\zeplus)\,\Vop\,\ket{\Psi^{(+)}}
= \ket{\Phi} + \Gzero(\zeplus)\,\sum_{(jk)} \Vij{jk}{}\,\ket{\Psi^{(+)}},
\end{equation}
其中 \(\ket{\Phi}\) 是入态。Faddeev 把右边的求和"拆分"成三段，
\begin{equation}
\label{eq:ch04-faddeev-decomp}
\ket{\Psi^{(+)}} \;=\; \sum_{i=1}^{3} \ket{\psifad{i}},
\quad
\ket{\psifad{i}} \;\equiv\; \Gzero\,\Vij{jk}{}\,\ket{\Psi^{(+)}},
\quad (i,j,k) \text{ 循环}.
\end{equation}
\(\ket{\psifad{i}}\) 称为\textbf{Faddeev 分量（Faddeev component）}。
注意 \(\Vij{jk}{}\) 只对 \((j,k)\) 这一对作用，所以 \(\ket{\psifad{i}}\) 是
"以粒子 \(i\) 为 spectator 的 12-3 树"那一族的"自然分量"。

把 \cref{eq:ch04-faddeev-decomp} 代回 \cref{eq:ch04-LS-full}：
\begin{equation}
\sum_i \ket{\psifad{i}}
= \ket{\Phi} + \Gzero \sum_{(jk)} \Vij{jk}{}\,\Big( \ket{\psifad{i}} + \ket{\psifad{j}} + \ket{\psifad{k}} \Big),
\end{equation}
其中 \((i, j, k)\) 仍循环排列。对每一个 \(i\) 单独取出来，并定义入态
\(\ket{\Phi^{(i)}}\) 为"对 \(\Vij{jk}{}\) 的散射前自由态投影"，得到
\textbf{Faddeev 方程组}：
\begin{equation}
\label{eq:ch04-faddeev-eqs}
\boxed{
\quad
\ket{\psifad{i}}
= \ket{\Phi^{(i)}}
+ \Gzero(\zeplus)\,\Vij{jk}{}\,\Big( \ket{\psifad{j}} + \ket{\psifad{k}} \Big),
\qquad i = 1, 2, 3.
\quad
}
\end{equation}
关键在于：右边\textbf{不再含 \(\Vij{jk}{}\,\ket{\psifad{i}}\)} 这一项——本来应
有的 \(\Gzero\,\Vij{jk}{}\,\ket{\psifad{i}}\) 项被 \cref{eq:ch04-faddeev-decomp}
中 \(\ket{\psifad{i}} = \Gzero \Vij{jk}{} \ket{\Psi}\) 的定义"吸进"了 LHS。
这正是 disconnected \(\delta\) 函数被消除的代数源头。

\paragraph{核的紧性。}
\cref{eq:ch04-faddeev-eqs} 的核是
\(K_{ij} = \Gzero\,\Vij{jk}{}\,\delta_{ij}^{(j \neq i)}\)（即 \(j = i\) 项不
出现）。它对每一个 \(i\)（行）只在 \(j \neq i\) 处有非零块，且每个非零块
都伴随着"换树"——即从 12-3 树 \(\to\) 23-1 树等的 Jacobi 旋转
（\cref{eq:ch03-tree-rotate}）。"换树"使 \(\bvec{q}^{(j)}\) 与 \(\bvec{q}^{(i)}\)
线性无关，于是三体动量积分对\textbf{所有 6 个动量分量}都是非平凡的，
不再有\(\delta\) 函数。Faddeev 1960 给出了核 \(K\) 在合适的能量上是 trace
class \(\Rightarrow\) 紧 \(\Rightarrow\) Fredholm 可逆的严格证明
（详见 Glöckle 1983 第 3 章）。

% -----------------------------------------
\subsection{用置换 \(\Pop\) 把三式压成一式}
\label{subsec:ch04-1plusP}

对全同粒子（如三个核子）三段 Faddeev 分量满足
\begin{equation}
\label{eq:ch04-Pop-on-psi}
\ket{\psifad{i}} = \Pop_i\,\ket{\psifad{1}},
\end{equation}
其中 \(\Pop_1 = \identity\)，\(\Pop_2 = P_{123}\)，\(\Pop_3 = P_{132} = P_{123}^2\)。
原因是循环对换把"以粒子 1 为 spectator 的分量"变到"以粒子 2 为 spectator 的分量"。
代入 \cref{eq:ch04-faddeev-eqs} 的 \(i = 1\) 式：
\begin{equation}
\ket{\psifad{1}}
= \ket{\Phi^{(1)}} + \Gzero \Vij{23}{} \big( \Pop_2 + \Pop_3 \big) \ket{\psifad{1}}
= \ket{\Phi^{(1)}} + \Gzero \Vij{23}{} \Pop \ket{\psifad{1}},
\end{equation}
其中
\begin{equation}
\label{eq:ch04-P-sum-def}
\Pop \;\equiv\; \Pop_2 + \Pop_3 \;=\; P_{123} + P_{132} \;=\; P_{123} + P_{123}^{-1}
\end{equation}
（注意：\(\Pop\) 这里是\textbf{两个循环对换之和}；它和 \cref{ch:03_jacobi_partial_wave}
里把 \(\identity\) 也算进去的 \(\Aop = \tfrac{1}{3}(\identity + P_{123} + P_{132})\) 差一个
\(\identity\) 项与归一因子）。代码里在 \texttt{P123\_sparse\_*} 数组中只
显式存 \(P_{123}\)；\(P_{132}\) 由 \(P_{123}\) 转置（行/列互换）得到。

\textbf{压缩后的 Faddeev 方程}：
\begin{equation}
\label{eq:ch04-faddeev-compact}
\boxed{
\quad
\ket{\psifad{}}
= \ket{\Phi}
+ \Gzero(\zeplus)\,\Vop\,\Pop\,\ket{\psifad{}}.
\quad
}
\end{equation}
（这里我们丢弃了下标 1，以"任一段"代表整族。\(\Vop\) 此处即 \(\Vij{23}{}\)
在 12-3 树上的形式；与 \cref{eq:ch04-LS-3body-naive} 的求和 \(\Vop\) 不同。）
这是单个的、紧 kernel 的、可数值化的方程。

\paragraph{用 \(\Tsub{}\) 重写。}
把 \(\Tsub{}\) 的 LS 关系（\cref{eq:ch04-tsub-def}）作用到 \cref{eq:ch04-faddeev-compact}：
\begin{equation}
\label{eq:ch04-faddeev-T-form}
\Tsub{}(\zeplus)\,\Pop\,\ket{\psifad{}}
= \Vop\,\big(\identity + \Gzero \Tsub{}\big)\,\Pop\,\ket{\psifad{}}.
\end{equation}
等价地——把 \cref{eq:ch04-faddeev-compact} 两边作用 \(\Vop\)、再套一次
LS 关系——可以推到\textbf{标准形式} \cite{Glockle1983}
\begin{equation}
\label{eq:ch04-T-from-faddeev}
\Top_{F}\,\ket{\Phi} = \Tsub{}\,\Pop \big( \identity + \Gzero \Top_F \big)\,\ket{\Phi},
\end{equation}
其中 \(\Top_F\) 是"Faddeev 形式的 T-算符"。这一形式与下面的 AGS \(U\) 算符
是等价的，相互之间只差代数重排。

% -----------------------------------------
\subsection{AGS 方程与 \(U\) 算符}
\label{subsec:ch04-AGS}

Alt-Grassberger-Sandhas（AGS）的工作把"三体 \(\to\) 三体"散射进一步组织成
"通道（channel）"间的转移振幅。一个"通道 \(\alpha\)"指的是在渐近自由态
区域 \(t \to \pm\infty\) 时保持束缚的子系统：
\begin{itemize}
  \item \(\alpha = 0\)：\textbf{断裂通道（breakup channel）}，三个粒子全自由。
  \item \(\alpha = 1\)：\textbf{1 + (23) 通道}，粒子 1 自由、(23) 形成束缚态
        （在核子物理里就是氘核 \(d\)）。
  \item \(\alpha = 2, 3\) 同理：\(2 + (31)\) 与 \(3 + (12)\)。
\end{itemize}
注意"通道 \(\alpha\)"\textbf{不是}\cref{ch:03_jacobi_partial_wave} 里的部分波
通道索引——前者是"渐近态结构"，后者是"角动量耦合标签"。Tic-tac 求解器
把两个含义都用 \(\alpha\) 当符号，难免歧义；本节中我们对前者用 \(\alpha, \beta\)，
对后者用 \(\PWchan, \PWchanp\)（与第 3 章一致）。

\paragraph{\(U_{\beta\alpha}\) 算符的定义。}
AGS 振幅 \(\Uampl{\beta}{\alpha}\)（"通道 \(\alpha \to \beta\) 的散射算符"）满足
\begin{equation}
\label{eq:ch04-AGS-master}
\boxed{
\quad
\Uampl{\beta}{\alpha} \;=\; \bar\delta_{\beta\alpha}\,\Gzero^{-1}
+ \sum_{\gamma=1}^{3} \bar\delta_{\beta\gamma}\,\Tsub{\gamma}\,\Gzero\,\Uampl{\gamma}{\alpha},
\quad
}
\end{equation}
其中 \(\bar\delta_{\beta\alpha} \equiv 1 - \delta_{\beta\alpha}\)。这是 9 个
方程组成的耦合方程组（\(\alpha, \beta = 1, 2, 3\)）。
"\(\bar\delta\)"的物理含义：当入射通道与出射通道相同时（如 \(d + p \to d + p\)，
即 \(\beta = \alpha\)），不再做"两体子内部散射"——它是渐近束缚态的内
部构造，已经在两体 T 算符之外了。

\paragraph{用 \(\Pop\) 把 9 式压成 1 式。}
对全同粒子，\cref{eq:ch04-Pop-on-psi} 的循环对称性让 \(\Tsub{2}, \Tsub{3}\) 都是
\(\Tsub{1}\) 经过 \(\Pop\) 共轭得到。沿用 \cref{subsec:ch04-1plusP} 的代数，
9 个 AGS 方程被压成单一的 \emph{对称形式}：
\begin{equation}
\label{eq:ch04-AGS-symm}
\Uop \;=\; \Pop\,\Gzero^{-1} + \Pop\,\Tsub{}\,\Gzero\,\Uop.
\end{equation}
其中 \(\Uop\) 是"对称化的"通道间振幅。两侧都左乘 \(\Vop\)：
\begin{equation}
\label{eq:ch04-AGS-V-form}
\Vop\,\Uop \;=\; \Vop\,\Pop\,\Gzero^{-1} + \Vop\,\Pop\,\Tsub{}\,\Gzero\,\Uop.
\end{equation}
利用 \(\Tsub{} = \Vop + \Vop\,\Gzero\,\Tsub{}\) 可以把 \(\Vop\,\Pop\,\Tsub{}\)
组合成 \((1 + \Vop \Gzero) \Tsub{} \Pop = \ldots\)（详细代数省略，见
Glöckle 1983 第 5 章）。最终把方程化为下面这条\textbf{Tic-tac 求解器实
际数值化的形式}：
\begin{equation}
\label{eq:ch04-faddeev-Mform}
\boxed{
\quad
\ket{M} \;=\; \Tsub{}\,\Pop \big( \identity + \Gzero \ket{M} \big),
\quad
\ket{M} \;\equiv\; \Uop\,\ket{\Phi},
\quad
}
\end{equation}
其中 \(\ket{\Phi}\) 是入态。把 \cref{eq:ch04-faddeev-Mform} 离散到部分波 + WP 网
格上得 \((\identity - K)\,U = K_0\) 的线性方程组，其中
\begin{equation}
\label{eq:ch04-K-kernel}
K = \Tsub{}\,\Pop\,\Gzero,
\qquad
K_0 = \Tsub{}\,\Pop\,\ket{\Phi}.
\end{equation}
这是 \cref{ch:11_faddeev_solver} 里 Padé 加速 Neumann 级数迭代的对象。

\paragraph{弹性通道 \(U_{0\beta}\) vs 断裂通道 \(U_{0(0)}\)。}
求解器内部统一构建上述 \(\Uop\)；从 \(\Uop\) 抽出物理观测量需要再把它投到
具体通道：
\begin{itemize}
  \item 弹性 \(d + p \to d + p\)：取入态 \(\ket{\Phi^{(d)}}\) 为氘核束缚态
        （由 SWP 系数 \(C_d\) 给出）+ 平面波 \(\ket{\bvec{q}_0}\)；
        然后 \(\Uampl{0}{\beta} = \langle d, \bvec{q}_0' | \Uop | d, \bvec{q}_0 \rangle\)。
        代码里这一步是 \texttt{post\_processing.cpp} 中通过
        \texttt{deuteron\_idx\_array} 与 \texttt{q\_com\_idx\_array}
        把 \(\Uop\) 三体矩阵元乘以 \(C_d\) 并对 \(p\) 维做求和的逻辑。
  \item 断裂 \(d + p \to N + N + N\)：取出态为三个自由波 \(\ket{\bvec{p}_0' \bvec{q}_0'}\)；
        \(\Uampl{0}{(0)}\)（"通道 0 \(\to\) 0"）的振幅是直接的
        \(\langle \bvec{p}_0' \bvec{q}_0' | \Uop | d, \bvec{q}_0 \rangle\)，
        在代码里对应 \texttt{U\_BU\_array}（"BU" = breakup）。
\end{itemize}
所以 \texttt{solve\_faddeev\_equations} 同时输出 \texttt{U\_array}（弹性的全 \((\PWchan, p, q)\)
矩阵元）与 \texttt{U\_BU\_array}（断裂的全矩阵元）。

% -----------------------------------------
\subsection{与 LS-form 的连接：\((1 + P_{12})\) 因子}
\label{subsec:ch04-LS-connect}

为了和"标准 LS 形式" \(\Top = \Vop + \Vop \Gzero \Top\) 直接对照，可以把
Faddeev/AGS 重写为
\begin{equation}
\label{eq:ch04-LS-link}
\Uop \;=\; (1 - P_{12}) \Tsub{}\,(1 + \Pop) \;+\; \ldots
\end{equation}
其中 \((1 - P_{12})\) 是反对称化投影因子（详见 \cref{eq:ch03-1plusP-projector}
的对称性等价物）。在等价表述中常见 "\((1 + P_{12})\,\Top\,\ldots\)"
的写法；它对应的是\textbf{对称}（玻色式）三体波函数；要用反对称（费米
式）则换成 \((1 - P_{12})\)。Tic-tac 求解器 \emph{隐式}处理这一因子——
反对称化通过 \cref{ch:03_jacobi_partial_wave} 中的 \(\ell + s + t\) 奇/偶
判据吸进了部分波态的归一化，再加上 \(\Pop = P_{123} + P_{123}^{-1}\) 自身
的对称性，最终给出与 Witała 等同的等价代数。

\textbf{硬性提醒}：在与外部 benchmark（如 Witała、Deltuva 提供的代码）对照
矩阵元时，请先确认对方有没有把 \((1 + P_{12})\) 显式拉出。如果对方有、
你这边没有，会差一个常数因子 2（或在反对称约定下差 \(-1\)，亦即整体取反）。

% =========================================
\section{离散化方案}
\label{sec:ch04-discretization}

\paragraph{从算符方程到矩阵方程。}
把 \cref{eq:ch04-faddeev-Mform} 投到 \(\pwstate{\PWchan}\) 与 WP 离散
\(\WPbin{i_p}, \WPbin{i_q}\) 的张量积基底：
\begin{equation}
\langle \PWchanp, i_p', i_q' | \Uop | \PWchan, i_p, i_q \rangle
\;\to\; (U)_{(\PWchanp, i_p', i_q'),(\PWchan, i_p, i_q)}.
\end{equation}
具体的 WP/SWP 离散化方法留到 \cref{ch:05_wpcd_math,ch:06_wp_swp_states}。
本节只交代"离散后的矩阵方程"长什么样。

\paragraph{密集化后的线性方程。}
求解器通过迭代解
\begin{equation}
\label{eq:ch04-discrete-master}
\big( \identity - K \big)\,U \;=\; K_0,
\qquad K = (C^T) (P_{123}) (V) (C),
\end{equation}
其中：
\begin{itemize}
  \item \(C \in \mathbb{R}^{(\Nalpha N_p^{\rm WP} N_q^{\rm WP}) \times (\ldots)}\)：
        Sturmian-Wave-Packet（SWP）展开系数，把 \(\Tsub{}\) 与 \(\Vop\) 的乘积凝
        缩成一个"两体 T 在 SWP 基下的列" —— 这正是为什么 \texttt{V\_WP\_*}
        与 \texttt{C\_WP\_*} 总是成对出现。
  \item \(P_{123}\)：\cref{ch:09_permutation_p123} 的稀疏算符，CSC 格式
        存储为 \texttt{P123\_sparse\_val\_array}, \texttt{P123\_sparse\_row\_array},
        \texttt{P123\_sparse\_col\_array\_csc}。
  \item \(V\)：两体势在 WP 基下的矩阵，分为未耦合 \(\Vop_{\rm unco}\)
        与耦合 \(\Vop_{\rm coup}\) 两块（对应 \texttt{V\_WP\_unco\_array} 与
        \texttt{V\_WP\_coup\_array}）。
  \item \(K_0\)：右端，等于 "\(K\) 作用在入态展开矢上"。
  \item 全维度 \(d = \Nalpha \cdot N_p^{\rm WP} \cdot N_q^{\rm WP}\)，
        对 \(\Nalpha = 365\), \(N_p^{\rm WP} = N_q^{\rm WP} = 30\) 给出
        \(d = 328\,500\)；矩阵密集存储是 \(d^2 \approx 10^{11}\) 个 double，
        约 800 GB——远超内存。代码因此\textbf{绝不显式构造 \(K\) 矩阵}，而是
        提供一个"\(K\) 作用在向量上"的列计算 callback（\texttt{CPVC\_col\_brute\_force}
        与其优化版），供 Padé/Neumann 迭代调用。
\end{itemize}

\paragraph{能量参数 \(\zeplus = E + i\eps\)。}
在物理上 \(\eps \to 0^+\) 给出散射极限。但 WP/SWP 离散化\textbf{自动}赋予
谱以一个有限分辨（即 WP bin 宽度）；这相当于在算子层面引入了一个
"组态平均"的有限 \(\eps\)。这就是为什么 Tic-tac 不需要在代码层显式选择
\(\eps\) 参数——它是 WP bin 宽度这一离散化选择的副产物。详见
\cref{ch:10_resolvent} 的 R + Q 分解。

\paragraph{3N-channel 块对角化。}
\(\Uop\) 关于 \(J^\pi\) 块对角，因此 \cref{eq:ch04-discrete-master} 是
按 \cref{ch:03_jacobi_partial_wave} 的 \texttt{chn\_3N\_idx\_array}
切分成若干独立块求解。每一块的维度是
\(N_{\alpha,\rm block} \cdot N_p^{\rm WP} \cdot N_q^{\rm WP}\)，
其中 \(N_{\alpha,\rm block}\) 不超过 63（最大块）。
不同 \(J^\pi\) 块在 OpenMP 外层并行。

% =========================================
\section{算法骨架}
\label{sec:ch04-algorithm}

\begin{algorithm}[H]
\caption{\texttt{solve\_faddeev\_equations} 的高层流程}
\label{alg:ch04-solve-faddeev}
\KwIn{$P_{123}$ 稀疏矩阵 (CSC)；$\Vop_{\rm unco}, \Vop_{\rm coup}$；
      SWP 系数 \texttt{swp\_states} (含 $C_{\rm unco}, C_{\rm coup}$)；
      free WP states \texttt{fwp\_states}；
      on-shell 索引数组 \texttt{chn\_os\_indexing}；
      部分波态空间 \texttt{pw\_states} (\cref{ch:03_jacobi_partial_wave});
      运行参数 \texttt{run\_parameters}；
      可选 3NF 模型 \texttt{tnf}}
\KwOut{弹性 $U$-数组 \texttt{U\_array}；断裂 $U_{\rm BU}$-数组；
       Green 函数 \texttt{G\_array}（仅在调试模式下保存）}
\BlankLine
\textbf{步骤 1（预构造）：} \;
$CT_{\rm RM} \leftarrow$ \texttt{create\_CT\_row\_maj\_3N\_pointer\_array}$(C_{\rm unco}, C_{\rm coup})$\;
$VC_{\rm CM} \leftarrow$ \texttt{create\_VC\_col\_maj\_3N\_pointer\_array}$(V_{\rm unco}, V_{\rm coup}, C_{\rm unco}, C_{\rm coup})$\;
\BlankLine
\textbf{步骤 2（3NF 上下文）：} \;
\If{\texttt{tnf} 非空}{
  组装 \texttt{tnf\_kernel\_context}（$p_{\rm WP}, q_{\rm WP}$ 边界 + scale)\;
}
\BlankLine
\textbf{步骤 3（按 3N-channel 块外层并行）：} \;
\For{每一 3N-channel block $c \in [0, N_{\rm chn,3N})$}{
  \For{每一 on-shell q-bin $i_{q_0} \in $ \texttt{q\_com\_idx\_array[c]}}{
    \For{每一 deuteron-state $i_d \in $ \texttt{deuteron\_idx\_array[c]}}{
      \BlankLine
      \textbf{步骤 3a：} 构造 $K_0$ 列（$=K \cdot \ket{\Phi^{(i_d)}}$）\;
      \textbf{步骤 3b：} Padé 加速 Neumann 迭代——
              每次迭代调用 \texttt{CPVC\_col\_brute\_force}
              （或 SIMD 优化版）逐列累乘 $K \cdot v$，
              一边累乘一边构造 Padé 表 $[n/n]$\;
      \textbf{步骤 3c：} 收敛判据基于相邻 Padé 阶差 $\nudif$；写入 \texttt{U\_array} 与 \texttt{U\_BU\_array}\;
    }
  }
}
\BlankLine
\textbf{步骤 4：} 释放 \texttt{CT\_RM\_array}, \texttt{VC\_CM\_array}\;
\Return\;
\end{algorithm}

\paragraph{复杂度分析。}
设 \(N_{\alpha,c}\) 是第 \(c\) 个 3N 块的部分波数，\(N_{\rm WP} = N_p^{\rm WP} N_q^{\rm WP}\)。
每次 \(K \cdot v\) 列计算的成本约
\(O(N_{\alpha,c}^2 \cdot N_{\rm WP} \cdot \mathrm{nnz}(P_{123})/N_{\alpha,c}) = O(N_{\alpha,c} \cdot N_{\rm WP} \cdot \mathrm{nnz}(P_{123}))\)。
Padé 收敛通常在 \(20\sim40\) 次迭代内完成。190 MeV 单 \(J^\pi\) 块、
\(N_{\rm WP} = 30 \times 30\) 在单线程下耗时约 30 秒—2 分钟，OpenMP 跨 8 核并行
后整套 dp 计算（10 块）约 5 分钟。

\paragraph{内存。}
\texttt{CT\_RM} 与 \texttt{VC\_CM} 是指针数组，指向\textbf{二体子块} \(V_{\rho} C_{\rho}\) 的固定 30×30 矩阵；
总内存 \(\sim O(\Nalpha^2)\) 个指针 + \(\sim O(\rho_{\max} \cdot N_p^2)\) 个 double，
对生产配置约 1 GB。

% =========================================
\section{代码落地}
\label{sec:ch04-code}

求解器入口位于 \texttt{src/core/faddeev\_solver/solve\_faddeev.h}（当前仓库）
和 \texttt{CPP/solve\_faddeev.h}（origin 仓库）。
本节给出函数签名以及辅助结构 \texttt{tnf\_kernel\_context}，
作为后续章节（特别是 \cref{ch:11_faddeev_solver}）的"接口契约"。

\paragraph{当前仓库（含 3NF）的接口签名。}
\lstinputlisting[
  style=cpp,
  caption={\texttt{solve\_faddeev\_equations} 函数签名 + 3NF 上下文（\texttt{Tic-tac/src/core/faddeev\_solver/solve\_faddeev.h:25--65}）},
  label={lst:ch04-faddeev-sig-current}
]{code_excerpts/snippets/ch04_solve_faddeev_signature.cpp}

\paragraph{tictac-origin（不含 3NF）的对应签名。}
\lstinputlisting[
  style=cpp,
  caption={origin 版本的 \texttt{solve\_faddeev\_equations} 签名（\texttt{tictac-origin/Tic-tac/CPP/solve\_faddeev.h:25--38}）},
  label={lst:ch04-faddeev-sig-origin}
]{code_excerpts/snippets/ch04_solve_faddeev_origin.cpp}

\paragraph{两签名的差异。}
\begin{itemize}
  \item \textbf{\texttt{fwp\_states} 参数}：当前仓库新增。这是 free
        Wave-Packet（FWP）边界数组（\(p_{\rm WP}, q_{\rm WP}\) 的 bin 边界），
        用于 3NF 计算需要的"WP 中心"参数。origin 仓库内不需要——它没有
        3NF 模块。
  \item \textbf{\texttt{tnf} 参数}：当前仓库的最后一个 default 实参
        \texttt{const three\_nucleon\_force\_model* tnf = nullptr}。
        origin 完全没有这一参数。当 \texttt{tnf == nullptr} 时（即不开 3NF），
        当前仓库的代码在所有 3NF 分支前都做一次空指针 check 后跳过——
        这是为了让"开 / 关 3NF"通过空对象（null object）模式实现，
        避免在调用方分裂出两条调用路径。
  \item \textbf{\texttt{tnf\_kernel\_context} 结构}：当前仓库新增的"3NF 数据
        包"。它把列计算热路径需要的所有 3NF 数据打包成一个 struct，
        避免函数签名指数级膨胀。注意 \texttt{w1\_scale} 字段——
        这是一个诊断旋钮，把整个 \(W^{(1)}\) 项乘以一个 scaling，
        用于 \cref{ch:17_3nf_numerical} 里的 sweep 测试。
\end{itemize}

\paragraph{每个参数的数学解释。}
\begin{description}
  \item[\texttt{U\_array}] \cref{eq:ch04-discrete-master} 的解，按
        \((\PWchan, i_p, i_q, i_{p_0}, i_{q_0})\) 索引；尺寸
        \(\Nalpha \cdot N_p \cdot N_q \cdot N_{p_0} \cdot N_{q_0}\)。
  \item[\texttt{U\_BU\_array}] 断裂通道矩阵元，按同样索引但 \(N_{p_0}\)
        换成"出射 \(p\) bin 数"。
  \item[\texttt{G\_array}] 三体 Green 函数 \(\Gthree(\zeplus)\) 在 SWP 基下的
        矩阵元；调试模式下打印用。
  \item[\texttt{P123\_sparse\_*}] CSC 格式 \(P_{123}\)。
        \texttt{val\_array}, \texttt{row\_array}, \texttt{col\_array\_csc}
        三件套；\texttt{P123\_sparse\_dim} 是稀疏元数。
        来自 \cref{ch:09_permutation_p123}。
  \item[\texttt{V\_WP\_unco\_array}, \texttt{V\_WP\_coup\_array}]
        二体势矩阵，未耦合 / 耦合分别存。来自 \cref{ch:08_potential_matrix}。
  \item[\texttt{swp\_states}] SWP 系数与维度，含 \texttt{C\_SWP\_unco/coup\_array},
        \texttt{Np\_WP}, \texttt{Nq\_WP}, \texttt{num\_2N\_unco/coup\_states}。
        来自 \cref{ch:06_wp_swp_states}。
  \item[\texttt{fwp\_states}] free WP 边界数组（仅当前仓库）。
  \item[\texttt{chn\_os\_indexing}] on-shell 索引：
        \texttt{q\_com\_idx\_array}（每个 \(\Tlab\) 在 q-bin 中映射到哪一格），
        \texttt{deuteron\_idx\_array}（氘核 SWP 索引），
        \texttt{num\_T\_lab}, \texttt{num\_deuteron\_states}。
  \item[\texttt{pw\_states}] 部分波态空间，由
        \cref{ch:03_jacobi_partial_wave} 的 \texttt{construct\_symmetric\_pw\_states}
        填好。
  \item[\texttt{file\_identification}] 输出文件标签字符串，控制
        \texttt{U\_PW\_elements\_*.txt} 命名。
  \item[\texttt{run\_parameters}] 运行参数（迭代上限、Padé 阶、收敛阈值、
        \texttt{w1\_scale} 等）。
  \item[\texttt{tnf}] 3NF 模型的可选注入；空指针 = 关 3NF。
\end{description}

\paragraph{辅助函数签名（同样在 .h 中，行 58--119）。}
本章不展开这些辅助函数的实现（详见 \cref{ch:11_faddeev_solver}），
但读者可以从签名直接看出三个观察：
\begin{itemize}
  \item \texttt{create\_CT\_row\_maj\_3N\_pointer\_array} 与
        \texttt{create\_VC\_col\_maj\_3N\_pointer\_array} 都\textbf{返回指
        针数组}（\texttt{double**}），表明 \(C^T\) 与 \(VC\) 在内存中按
        2N 子矩阵分块；指针数组的索引则按 \((\PWchan, \PWchanp)\) 二维。
  \item \texttt{PVC\_col\_brute\_force} 与 \texttt{CPVC\_col\_brute\_force}
        是\textbf{列计算} routine——给定 \((i_\PWchan, i_p, i_q)\) 列号，
        计算 \(K \cdot e_{(\PWchan, i_p, i_q)}\) 这一列，返回到 \texttt{col\_array}。
        这是为了避免显式构造 \(K\) 矩阵。
  \item \texttt{*\_col\_calc\_test} 是\textbf{比对工具}：把"暴力构造完整
        矩阵再乘"与"直接调用 col 计算"两条路径的输出对比，确保正确。
        默认关闭（\texttt{test\_PVC\_col\_routine = false}），仅 benchmark 时打开。
\end{itemize}

% =========================================
\section{数字闭环}
\label{sec:ch04-numerics}

\begin{numericalbox}[label={box:ch04-numerical-closure}]
\textbf{场景}：固定 \(J^\pi = 1/2^+\) 块，\(J_{2N\max} = 3\)，
\(2J_{3N\max} = 1\)，\(N_p^{\rm WP} = N_q^{\rm WP} = 30\)。
求该块上 AGS 矩阵 \cref{eq:ch04-discrete-master} 的维度与稀疏度。

\paragraph{(a) 部分波数。}
由 \cref{box:ch03-numerical-closure} (c) 已知 \(\Nalpha(J^\pi = 1/2^+) = 27\)。

\paragraph{(b) 单块矩阵维度。}
\begin{equation}
d_{\rm block} = \Nalpha \cdot N_p^{\rm WP} \cdot N_q^{\rm WP}
= 27 \cdot 30 \cdot 30 = 24\,300.
\end{equation}
若密集存储 \(K\)（双精度），需要 \(24\,300^2 \times 8\,\text{B} \approx 4.7\,\text{GB}\)——单块也已不便密存。

\paragraph{(c) 稀疏度估算。}
\(P_{123}\) 在通道层是\textbf{密}的（约 \(\Nalpha^2 = 729\) 个非零通道对），
但每个通道对内部沿 \((i_p, i_q)\) 维只在 "Jacobi 旋转后落到的 1\(\sim\)8 个 bin"
里非零。
经验上 \(P_{123}\) 的稀疏度 \(\rho_P \sim 1\%\)，
因此 \texttt{P123\_sparse\_dim} \(\approx 0.01 \times d_{\rm block}^2 \approx 6 \times 10^6\) 元，
约 100 MB（每元 = 1 个 double + 1 个 size\_t \(\approx 16\) B）。

\paragraph{(d) 复现命令。}
\begin{lstlisting}[style=config]
# 启用 P123 稀疏度的诊断打印
cd Tic-tac/CPP
./run J_2N_max=3 two_J_3N_max=1 \
      Np_WP=30 Nq_WP=30 \
      potential_model=N2LOopt \
      Tlab=190.0 \
      output_folder=Output/ch04_dim_check \
      reuse_p123=false \
    2>&1 | tee /tmp/ch04_dim.log

# 读 P123 实际稀疏元数
grep -i "P123_sparse_dim\|nnz" /tmp/ch04_dim.log
\end{lstlisting}
预期 \texttt{P123\_sparse\_dim} 在 \(10^6\) 量级（具体值依赖
3N-channel 选择与 \texttt{chebyshev\_s, chebyshev\_t} 网格参数）。

\end{numericalbox}

% =========================================
\section{接口与依赖}
\label{sec:ch04-interface}

本章定义的 \(\Uop\) 算符与 AGS 矩阵 \(K\) 是后续\textbf{所有}章节的核心。
以下是依赖拓扑：

\paragraph{上游依赖（提供给本章 \(\Uop\) 求解所需的输入）。}
\begin{itemize}
  \item \cref{ch:03_jacobi_partial_wave}：\(\Nalpha\) 与 9 个通道索引数组。
  \item \cref{ch:05_wpcd_math,ch:06_wp_swp_states}：\(\WPbin{i}\) 基底与
        \(C\) 矩阵。
  \item \cref{ch:08_potential_matrix}：\(V\) 矩阵。
  \item \cref{ch:09_permutation_p123}：\(P_{123}\) 稀疏矩阵。
  \item \cref{ch:10_resolvent}：\(\Gzero\) 在 WP 基下的对角元（即"通道
        Resolvent"）。
\end{itemize}

\paragraph{下游消费（消费 \(\Uop\) 输出的章节）。}
\begin{itemize}
  \item \cref{ch:11_faddeev_solver}：本章 \cref{eq:ch04-faddeev-Mform}
        的 Padé/Neumann 数值实现；输出 \texttt{U\_array}, \texttt{U\_BU\_array}。
  \item \cref{ch:12_observables_basics}：从 \texttt{U\_array} 提取弹性散
        射振幅 \(M_{m_d', m_d}(\theta)\) 与微分截面 \(\diffXS\)。
  \item \cref{ch:13_polarization_observables}：从同一 \(U\) 提取
        \(\iTeleven, \Ttwentyzero, \Ttwentyone, \Ttwentytwo\)。
  \item \cref{ch:14_miller_gate1}：用 \(U\) 在小角度处提取相移
        （phase shift），即 Miller Gate 1 程式。
  \item \cref{ch:15_3nf_physics,ch:16_3nf_pw_projection,ch:17_3nf_numerical}：
        通过 \texttt{tnf} 注入修改 \(K\)；再次走本章流程。
\end{itemize}

\paragraph{依赖图。}
\begin{figure}[h]
\centering
\begin{tikzpicture}[
  node distance=4mm and 8mm,
  every node/.style={draw, rounded corners=2pt, font=\footnotesize,
                     inner sep=3pt, align=center},
  blk/.style={draw=black, fill=blue!8},
  outblk/.style={draw=red!50!black, fill=red!5},
  arr/.style={-Stealth, thick}
]
  \node[blk] (ch3) {Ch 3\\$\Nalpha,\PWchan$};
  \node[blk, right=of ch3] (ch6) {Ch 6\\$C$ (SWP)};
  \node[blk, right=of ch6] (ch8) {Ch 8\\$V$ matrix};
  \node[blk, right=of ch8] (ch9) {Ch 9\\$P_{123}$};
  \node[blk, right=of ch9] (ch10) {Ch 10\\$\Gzero$ (R+Q)};
  \node[draw, fill=yellow!30, below=12mm of ch8, minimum width=5cm] (ch4)
    {\textbf{Ch 4 (本章)}\\AGS / $\Uop$ 算符\\
     \texttt{solve\_faddeev\_equations}};
  \node[blk, below=of ch4] (ch11) {Ch 11\\Faddeev 求解器\\(Padé+Neumann)};
  \node[outblk, below=10mm of ch11, xshift=-3.5cm] (ch12) {Ch 12 截面};
  \node[outblk, below=10mm of ch11, xshift=-1cm]   (ch13) {Ch 13 极化};
  \node[outblk, below=10mm of ch11, xshift=1.5cm]  (ch14) {Ch 14 相移};
  \node[outblk, below=10mm of ch11, xshift=4cm]    (ch17) {Ch 15-17 3NF};
  \draw[arr] (ch3) -- (ch4);
  \draw[arr] (ch6) -- (ch4);
  \draw[arr] (ch8) -- (ch4);
  \draw[arr] (ch9) -- (ch4);
  \draw[arr] (ch10) -- (ch4);
  \draw[arr] (ch4) -- (ch11);
  \draw[arr] (ch11) -- (ch12);
  \draw[arr] (ch11) -- (ch13);
  \draw[arr] (ch11) -- (ch14);
  \draw[arr, dashed] (ch17) -- (ch4) node[midway, right, font=\scriptsize] {\texttt{tnf} 注入};
\end{tikzpicture}
\caption{第 4 章 AGS / \(\Uop\) 算符的上下游依赖与下游消费拓扑。}
\label{fig:ch04-callgraph}
\end{figure}

% =========================================
\section{常见陷阱}
\label{sec:ch04-pitfalls}

\begin{pitfallbox}
\textbf{陷阱 1：\(\Tampl{\beta}{\alpha}\) vs \(\Uampl{\beta}{\alpha}\) 的符号约定。}

不同文献对"通道间转移振幅"使用 \(T_{\beta\alpha}\) 与 \(U_{\beta\alpha}\) 两套
记号，二者相差一个常数因子（通常 \(\pm 1\) 或 \(\pm i\)，依归一约定）。
Tic-tac 内部统一用 \(U\)（即 AGS 表述）。在与 Witała 1985 PRC 的相
移表对照时请\textbf{先确认对方用 T 还是 U}：差一个 \(-1\) 会让 Ay 倒符号
但 \(d\sigma/d\Omega\) 不变（因为后者只取 \(|U|^2\)）。

\textbf{自检}：用 LO\_internal 势在 190 MeV 跑 Ay 应得\textbf{正}峰值（约
+0.4）；若得到 \(-0.4\)，多半是这一符号搞错。
\end{pitfallbox}

\begin{pitfallbox}
\textbf{陷阱 2：复能量 \(\zeplus = E + i\eps\) 的位置。}

\cref{eq:ch04-faddeev-Mform} 与 \cref{eq:ch04-AGS-symm} 都隐式依赖
\(\zeplus = E + i\eps\) 的虚部 \(\eps\)。WP 离散化\textbf{自动}给出有效
\(\eps_{\rm eff} \sim \Delta q_{\rm WP}^2 / (2 \mu)\)（即 bin 宽度的能量分
辨）。
\textbf{陷阱所在}：如果你尝试在求解器外手动加 \(i\eps\)（如修改 \texttt{make\_resolvent.cpp}
的实部），会与 WP 的内置正则化叠加，得到双重正则化、共振峰被人为加宽。
正确做法是\textbf{增大 \(N_q^{\rm WP}\)} 而不是显式加 \(\eps\)。
\end{pitfallbox}

\begin{pitfallbox}
\textbf{陷阱 3：与 LS 形式互转时遗漏 \((1 + P_{12})\) 因子。}

\cref{subsec:ch04-LS-connect} 已经讲过：Tic-tac 把 \((1 - P_{12})\)
\textbf{隐式}吸进部分波态归一化，外部 benchmark（如 Deltuva 的代码）通
常\textbf{显式}写 \((1 + P_{12})\)。直接用矩阵元做对照时差一个因子
2（玻色约定）或一个负号（费米约定）。

\textbf{对照建议}：取一个在两个代码里都明确给出的"参考量"——
通常用弹性总截面 \(\sigma_{\rm el}\)（积分过 \(\theta\) 后的 \(d\sigma/d\Omega\)）做闭环，
而不是裸矩阵元。
\end{pitfallbox}

\begin{pitfallbox}
\textbf{陷阱 4：\texttt{q\_com\_idx\_array} 与 \texttt{deuteron\_idx\_array} 的"隐式
长度对偶"。}

\texttt{chn\_os\_indexing} 包含 \texttt{num\_T\_lab} 个 q-bin 索引（即
"输入要求的能量在 q-WP 网格上落到第几个 bin"）和
\texttt{num\_deuteron\_states} 个氘核 SWP 索引。这两个数组\textbf{独立长度}，
但被嵌套循环展开为
\texttt{num\_T\_lab \(\times\) num\_deuteron\_states} 次实际 \(K_0\) 列计算。

\textbf{陷阱所在}：
若 \texttt{deuteron\_idx\_array} 为空（如某些模型不支持氘核束缚），
求解器会跳过整个 elastic 通道但仍保留 BU 数组分配，
最终输出 \texttt{U\_array} 全为零的 NaN。
日志里只会打印 "0 deuteron states"，容易漏看。
\textbf{自检}：在 PR 修改任何 SWP 构造代码后，
跑一遍 \texttt{python tests/test\_190mev\_data\_pipeline.py}——
它会显式检查氘核态数量是否大于零。
\end{pitfallbox}

\begin{pitfallbox}
\textbf{陷阱 5：\texttt{tnf == nullptr} 与 \texttt{tnf->enabled() == false} 的差别。}

当前仓库设计上有两套"关 3NF"的方式：
(a) 把 \texttt{tnf} 直接传 \texttt{nullptr}；
(b) 传一个非空但内部 \texttt{enabled() == false} 的 stub 模型。
方式 (a) 完全跳过任何 3NF 相关分支；方式 (b) 走入分支但每个内部循环再做一次 \texttt{enabled()} 检查。

\textbf{陷阱所在}：
两种调用路径在某些数值边界（如 P123 缓存 hash key 计算）会得到\textbf{不同}的中间状态——
即缓存命中与不命中——尽管最终物理结果一致。
\textbf{建议}：所有生产 / sweep 跑统一走方式 (a)；只在调试 3NF 注入路径时才用 (b)。
否则可能复现不出来上次的 \texttt{P123\_cache\_*.h5} 缓存命中。
\end{pitfallbox}

\begin{pitfallbox}
\textbf{陷阱 6：\texttt{file\_identification} 字符串与缓存文件名。}

\texttt{file\_identification} 是 \texttt{solve\_faddeev\_equations} 的字符串
参数，最终拼到 \texttt{U\_PW\_elements\_*.txt} 与 \texttt{neumann\_terms\_*.txt}
文件名里。如果你跑一个 sweep（如多个 \texttt{w1\_scale} 值）但忘了改这个字段，
后续跑\textbf{会覆盖前一次的输出}。
\textbf{建议}：在 sweep 脚本里把所有可调参数（\(\Tlab\), \(N_p, N_q\), 势模
型, \texttt{w1\_scale}）都拼进 \texttt{file\_identification}，并在 commit 信
息里写清"本次 sweep 输出文件名规则"。
\end{pitfallbox}
